\documentclass[11pt]{article}
\usepackage[margin=1in]{geometry}
\usepackage{amsmath,amssymb,mathtools}
\usepackage{fontspec,unicode-math}
\setmainfont{Latin Modern Roman}
\setsansfont{Latin Modern Sans}
\setmathfont{Latin Modern Math}
\usepackage{xurl,hyperref}
\hypersetup{colorlinks=true,urlcolor=blue,linkcolor=blue}
\newcommand{\sref}[2]{\hyperlink{source:#1:#2}{[S#2]}}
\newcommand{\cataloguescope}[1]{\paragraph{Recorded mathematical question.}#1}
\begin{document}
% UnsolvedMath ID 9400177; source catalogue code AMR-093-0177
\setcounter{section}{9400176}
\section{Pollock's octahedral-number conjecture}\label{9400177}
{\small\sffamily UnsolvedMath ID 9400177 · Number Theory}
\subsection{Definitions and mathematical statement}

\paragraph{Shared notation.}$\mathbb N=\{1,2,\ldots\}$, and $\mathbb Z,\mathbb Q,\mathbb R,\mathbb C$ denote the integers, rationals, reals and complex numbers. Any other conventions are specified locally.


For $m\in\mathbb Z_{\ge0}$ define the octahedral number $O_m=(2m^3+m)/3$. These are nonnegative integers, with $O_0=0$. Thus a sum of at most seven positive octahedral numbers can be padded with zeros to a sum of exactly seven members of this sequence.
\paragraph{Statement recorded in the supplied source.}
Is it true that for every $n\in\mathbb N$ there exist $m_1,\ldots,m_7\in\mathbb Z_{\ge0}$ such that $n=\sum_{j=1}^7 O_{m_j}$? \sref{9400177}{2}
\subsection{Short English statement}
Can every positive integer be written as the sum of at most seven octahedral numbers?
\subsection{Sources}
\begin{description}
\item[\hypertarget{source:9400177:1}{S1}] ULAM.AI and the UnsolvedMath Contributors, \emph{UnsolvedMath: A Curated Collection of Open Mathematics Problems}, record 9400177 (AMR-093-0177). CC BY 4.0. \url{https://huggingface.co/datasets/ulamai/UnsolvedMath}.
\item[\hypertarget{source:9400177:2}{S2}] Zarathustra Brady, \emph{Sums of seven octahedral numbers}, arXiv:1509.04316v2 (2015), Section 1, Pollock conjecture and definition of octahedral numbers, pp. 1--2. \url{https://arxiv.org/pdf/1509.04316}.
\end{description}
\label{end:9400177}
\end{document}
